\section{Related Work}
% \textbf{Knowledge Distillation}~\cite{buciluǎ2006model, hinton2015distilling} is a technique that trains a small model to approximate the output distribution of a large model. With recent advances in LLMs, various distillation approaches have been proposed for auto-regressive models, including matching token-level distribution~\cite{sanh2019distilbert}, teacher-generated sequences~\cite{kim2016sequence}, or attention scores~\cite{wang2020minilm}. In addition, several studies have explored the alternative divergence objectives~\cite{wen2023f, ko2024distillm, shing2025taid} to overcome the limitations of forward and reverse KL and to improve the stability of distribution matching. Also, to address exposure bias caused by the mismatch between teacher-generated training data and self-generated sequences at inference time \cite{bengio2015scheduled, ranzato2015sequence}, several on-policy knowledge distillation methods have been proposed.
% These include methods that combine off-policy and on-policy data \cite{lin2020autoregressive}, techniques focusing on reverse KL divergence over student-generated contexts \cite{gu2023minillm, lu2025onpolicydistillation}, an on-policy distillation pipeline with multiple objectives \cite{agarwal2024policy}, and an approach based on interleaved sampling \cite{xu2024speculative}. In this paper, we build upon on-policy distillation~\cite{lu2025onpolicydistillation} by leveraging dense teacher supervision and effectively transferring the teacher’s uncertainty based on its entropy.

\textbf{Knowledge Distillation}~\cite{bucilua2006model, hinton2015distilling} trains a smaller model to approximate a larger model's output distribution. For auto-regressive models, various approaches have been proposed, including matching token-level distributions~\cite{sanh2019distilbert}, teacher-generated sequences~\cite{kim2016sequence}, attention scores~\cite{wang2020minilm}, and alternative divergence objectives~\cite{wen2023f, ko2024distillm, shing2025taid} to overcome limitations of forward and reverse KL. To address exposure bias from the mismatch between teacher-generated training data and self-generated sequences at inference~\cite{bengio2015scheduled, ranzato2015sequence}, several on-policy methods have been proposed, including combining off-policy and on-policy data~\cite{lin2020autoregressive, agarwal2024policy}, reverse KL over student-generated contexts~\cite{gu2024minillm, lu2025onpolicydistillation}, and interleaved sampling~\cite{xu2024speculative}. 
In addition, approaches that adapt the divergence objective based on teacher-student discrepancy measures, such as entropy gaps~\cite{amara2022bd}, head-tail mismatch~\cite{wu2025rethinking}, and vocabulary-level discrepancy~\cite{jung2025todi}, have also been explored. Our work specifically focuses on the instability and diversity degradation caused by reverse KL in high-entropy regions under on-policy distillation, selectively applying forward KL based on teacher uncertainty.

% \textbf{Reasoning Abilities of Language Models} have been improved by various approaches, including prompting-based methods that enhance reasoning without further post-training~\cite{wei2022chain, yao2023tree}, test-time scaling techniques that increase computation during inference~\cite{muennighoff2025s1, snell2024scaling}, and supervised fine-tuning approaches that leverage outputs generated by larger models~\cite{guha2025openthoughts, he2025deepmath, guo2025deepseek}.
\textbf{Reasoning Abilities of Language Models} have advanced through prompting~\cite{wei2022chain, yao2023tree}, test-time scaling~\cite{muennighoff2025s1, snell2024scaling}, and distillation from larger models~\cite{guha2025openthoughts, he2025deepmath, guo2025deepseek}.
Recently, RL methods have received particular attention, including methods that optimize verifiable rewards~\cite{shao2024deepseekmath, yu2025dapo, liu2503understanding} or apply step-level supervision using process reward models~\cite{uesato2022solving, lightman2023let}. In addition, \citet{cheng2025reasoning, wang2025beyond} have proposed entropy-based methods to encourage exploration during training. 
Instead of relying on intrinsic entropy for exploration, our method uses teacher-guided KL selection to transfer uncertainty and distributional structure.
% Instead of encouraging exploration via intrinsic entropy, our method uses teacher-guided KL selection to transfer uncertainty and distributional structure.